\section{Games on graphs}\label{sec:games}

In a mean-payoff game, two players choose a walk in a directed graph.
Max wants the walk to earn a high average reward and Min wants it to
earn a low one. We compute the value from every starting state and
optimal policies for both players. The algorithm is strongly polynomial
when each action determines its outcome. If an action draws a random
outcome, the step bound also counts the bits specifying its
probabilities. The proof combines two earlier results. A fixed-point
computation proposes actions for the two players, and linear programming
determines what each proposed policy guarantees against a best response.

\paragraph{Arena and payoffs.}
The arena is finite and explicit. Each of its \(n\ge1\) states belongs
to Max or Min, whose turn it is to choose an action there. Every state
has at least one action. An action \(a\) has a nonempty finite set
\(E_a\) of outcomes. Outcome \(e\) leads to state \(j(e)\) and pays a
real reward \(r_e\). Conditional on the entire history and the chosen
action, \(e\) has probability \(p_e\). Each action belongs to a single
state and each outcome to a single action. Players observe all actions
and outcomes and may randomize as a function of the entire history.
A \emph{policy} is a pure positional strategy, specifying one action
at each state owned by that player.

For a play with outcomes \(e_1,e_2,\ldots\), the ordinary mean payoff is
\[
    \liminf_{N\to\infty}\frac1N\sum_{k=1}^N r_{e_k}.
\]
We also allow \emph{priority mean payoff} \cite{GZ}. Each outcome has
a weight \(w_e>0\) and a positive integer priority \(q_e\). Let
\(q=\liminf_k q_{e_k}\), the least priority visited infinitely often.
The payoff is
\[
  \liminf_{N\to\infty}
  \frac{\sum_{k\le N:\,q_{e_k}=q} w_{e_k}r_{e_k}}
       {\sum_{k\le N:\,q_{e_k}=q} w_{e_k}} .
\]
The denominator is eventually positive. Equal priorities and unit weights
give ordinary mean payoff. Priorities are discrete input and can be
replaced by their ranks \(1,\ldots,Q\), with \(Q\) at most the number
of outcomes. In a \emph{discounted} game, each outcome has a discount
factor \(0\le\beta_e<1\). Its payoff is
\[
   \sum_{k\ge1} r_{e_k}\prod_{l<k}\beta_{e_l}.
\]
For each of these payoffs Max maximizes its expectation and Min
minimizes it.

\begin{theorem}[Values and optimal policies]\label{thm:games}
For discounted and priority mean-payoff games on finite turn-based
arenas of perfect information, there is a deterministic algorithm that
computes the exact value vector \(v\) and a pair of policies. From
every initial state \(i\), the Max policy guarantees expected payoff
at least \(v_i\) and the Min policy guarantees at most \(v_i\) against
every strategy of the opposing player.
\begin{itemize}
 \item In a two-player game with binary rational transition
 probabilities, the step bound is polynomial in the size of the arena,
 including its discrete data and the transition probability bit
 lengths. Rewards, weights, and discount factors may be arbitrary real
 scalars.
 \item In a Markov decision process, where only one player has choices,
 the step bound is polynomial in the size of the arena and its discrete
 data even for arbitrary real transition probabilities supplied as scalars.
\end{itemize}
On rational input all intermediate numbers and outputs have bit lengths
polynomial in the complete input length. In particular, both cases have
polynomial bit complexity on binary rational input.
\end{theorem}

In a deterministic arena there is one outcome per action, so the
two-player algorithm is strongly polynomial, including for arbitrary
real rewards. Against pure strategies the policy guarantees hold for
the resulting individual play. For both mean-payoff objectives the
same values and policies apply when the liminf is replaced by a limsup.

\paragraph{Reachability and parity.}
Making a target absorbing and assigning reward one there and zero
elsewhere turns the mean payoff of each play into the indicator of
reaching the target. Theorem~\ref{thm:games} therefore gives a
polynomial-time algorithm for simple stochastic games, whose threshold
problem Condon placed in NP \(\cap\) coNP \cite{Cond}. For reductions
among simple stochastic games, discounted games and mean-payoff games,
see Zwick--Paterson \cite{ZP} and Andersson--Miltersen \cite{AM}.

In a parity game Max wins when the least priority occurring infinitely
often is even. Set \(r_e=\mathbf 1_{\{q_e\ {\rm even}\}}\) before
relabeling the priorities. The priority mean payoff of each play then
equals its win indicator \cite{GZ}, so the algorithm computes the
optimal win probability even with binary rational transitions.
The deterministic parity problem also reduces to ordinary deterministic
mean payoff with exponentially large rewards \cite{Jur}. Earlier
quasipolynomial algorithms for deterministic parity games include the
breakthrough of Calude, Jain, Khoussainov, Li and Stephan \cite{Cal}
and Parys's recursive approach \cite{Parys}. The distinct
\emph{mean-payoff parity} objective imposes a parity condition on a
mean-payoff game \cite{OAIMe}.

\paragraph{Earlier algorithms.}
Shapley \cite{Shap} introduced discounted stochastic games and the
contraction argument, allowing simultaneous moves. Gillette \cite{Gillet}
studied long-run average reward. Positional optimality for
perfect-information games has a long history, including
Liggett--Lippman \cite{LL} and, for deterministic games,
Ehrenfeucht--Mycielski \cite{EM}.

Pseudopolynomial algorithms for integer rewards use finite horizons
\cite{ZP} or an energy objective with initial credit \cite{BC}.
See also \cite{GKK,Puri} and algorithms parameterized by the number
of chance positions \cite{BEGM,AGKS}. Randomized strategy iteration
draws on subexponential simplex methods \cite{Kalai,MSW,HZ}. It has
been applied to simple stochastic games \cite{Lud} and mean-payoff
games \cite{BSV,BV}, with an LP-type interpretation \cite{Hal} and
improved expected bounds \cite{Uz}. Complementarity formulations are
available for discounted and simple stochastic games
\cite{JSav,HIJ,GR}, although their matrices can fail the monotonicity
hypothesis of Section~\ref{sec:quadratic} \cite[Theorem~23]{HIJ}.

Ohlmann \cite{Ohl} gave a symmetric recursive algorithm for
deterministic mean-payoff games. Recent deterministic and randomized
algorithms give quasipolynomial bounds for binary-encoded integer
rewards \cite{OAIDe,OAIRa}, and a deterministic quasipolynomial
bound is known for turn-based stochastic games \cite{OAITu}. For
connections with contractions see \cite{SF}. Truffet's proposals
use max-atoms and tropical optimization \cite{Tru1,Tru2}. See
\cite[Section~1 and Appendix~A]{OAIRa} for discussion.

\paragraph{Proof of Theorem~\ref{thm:games}.}
We first solve the discounted fixed-point equation when its slopes
have short binary encodings. Taking discounts to one then yields
mean-payoff values and policies by Blackwell optimality. To remove
the dependence on discount bits, we use LP to evaluate a fixed policy
with arbitrary real data. These one-player evaluations give lower and
upper guarantees on the two-player value. A clipped fixed-point
computation halves the gap between the guarantees, and a bound on
the possible scales of a positive gap gives exact termination.

\subsection{Discounted values and the mean-payoff limit}
\label{sec:shortslopes}

At a state \(i\), write \(\mathrm{ext}_i\) for the maximum over actions
if Max owns \(i\) and for the minimum if Min owns it. The Bellman
(Shapley) operator for a discounted game is
\begin{equation}\label{eq:T}
 T_i(x)=\mathop{\mathrm{ext}_i}_{a\text{ at }i}
        \sum_{e\in E_a}p_e(r_e+\beta_e x_{j(e)}).
\end{equation}
It preserves the coordinate order and contracts the sup norm by a
factor at most \(\bar\beta=\max_e\beta_e<1\). Suppose \(x\le T(x)\).
At each Max state choose an action attaining the maximum. At a Min
state every action has score at least \(x_i\). Under this Max policy,
iteration of the one-step inequality gives expected discounted reward
through \(N\) steps, including discounted terminal value \(x\), at
least \(x_i\) from start \(i\). The terminal contribution has magnitude
at most \(\bar\beta^N\|x\|_\infty\), so the same guarantee holds for
infinite play. Reversing the inequality gives Min's upper guarantee.
At the fixed point these bounds coincide. An extremizing action in
each state gives an optimal policy for its owner.

Suppose the discount factors and probabilities are given as binary
rationals. Scale the rewards by \(1/S\), where
\[
 S=\max\{1,\max_e |r_e|/(1-\beta_e)\}.
\]
Every term \(r_e/S+\beta_e x_{j(e)}\) then lies in \([-1,1]\) on the
unit cube. The circuit for~\eqref{eq:T} has binary rational slopes
and arbitrary real offsets, as required by Proposition~\ref{prop:asym}.
Multiplying its fixed point by \(S\) recovers the value vector.
The extremizing actions survive the scaling. The arithmetic uses
only a constant number of layers of scalar formulas and subroutine calls.

The same contraction estimate gives an approximation result for
discounted games \emph{with simultaneous moves}. Here the players
choose independently randomized actions at each state, and the
outcome distribution depends on the pair of actions. All history is
public. Shapley's operator replaces the extremum in~\eqref{eq:T} by
the value of the one-stage matrix game \cite{Shap}. Optimal mixed
actions at its fixed point give the same finite-horizon guarantee.
With the scaling above, the map preserves the cube and contracts it by
\(\bar\beta\), since an entrywise perturbation of size \(\epsilon\)
changes a matrix-game value by at most \(\epsilon\). For binary
rational input of length \(L\), \(1-\bar\beta\ge 2^{-L}\). Each
evaluation can use bit-polynomial LP. Theorem~\ref{thm:contraction}
therefore approximates the value vector to absolute error \(2^{-k}\)
in bit time polynomial in \(L\) and \(k\). Undoing the scaling costs
only polynomially many extra bits of accuracy. See
\cite{EPRY,BFGMS,FH} for related fixed-point methods.

\paragraph{From discounts to mean payoff.}
In a turn-based priority game, introduce the parameter \(t>0\) and set
\begin{equation}\label{eq:gz}
   \beta_e=1-w_e t^{q_e},\qquad
   \widetilde r_e=(1-\beta_e)r_e .
\end{equation}
Use \(q_e\) after relabeling. For small \(t\), the resulting discounted
payoff is the expected reward in a process that stops after outcome
\(e\) with probability \(1-\beta_e\) and pays \(r_e\) upon stopping
\cite{GZ}. Lower priorities have larger stopping probabilities as
\(t\) tends to zero.

For binary rational probabilities and weights, the slopes in the
discounted operator are polynomials with binary rational coefficients.
The scaling \(S=\max\{1,\max_e|r_e|\}\) works for every small \(t>0\).
Proposition~\ref{prop:asym} supplies the discounted values as rational
functions of \(t\) and a pair of \emph{Blackwell optimal} policies,
optimal for all sufficiently small \(t>0\). For background on
Blackwell optimality see \cite{Bla,AHM,GZDet,GZMD}. Determinant signs
have also been used to establish stabilization as discounts tend to one
\cite{AM,GGK}. See also \cite[Section~2.1]{OAITu}.

Gimbert and Zielonka's Theorem~6 \cite{GZ} transfers these policies
to priority mean payoff and identifies the values with the limits of
the discounted values. Their theorem uses rewards, priorities and
weights on states. To apply it here, duplicate each state to record
its incoming outcome. Give each original state an additional starting
copy with fixed dummy data. From a specified starting copy the observed
histories in the two games determine one another, and an original
policy lifts by ignoring the incoming outcome. The initial state
datum adds a fixed first discounted payoff and multiplies the
remaining payoff by a positive factor tending to one. The additive
term tends to zero. Thus the lifted policies are Blackwell optimal,
and the extra initial datum does not affect priority mean payoff.
Applying the theorem with parameter \(1-t\) gives its limsup
conclusion. Exchanging players and negating the rewards gives the
liminf conclusion with the same policies.

The transfer holds for arbitrary real game data. The encoding
assumptions above enter the computation of the Blackwell policies.
Once the value functions are known, their limits are obtained from
their lowest nonzero coefficients. The values are bounded, and an
identically zero numerator gives value zero.

\subsection{The one-player case}\label{sec:mdp}

In a Markov decision process (MDP) only one player chooses among
actions. After negating rewards if necessary, take this player to be
Max. We allow arbitrary real probabilities. More generally, we allow
rewards and discounts to be polynomials in \(t\) of explicitly given
degree, provided that the discounts belong to \([0,1)\) for all
sufficiently small \(t>0\). We will compute a single optimal policy
and rational value functions valid throughout such an interval.

Dynamic programming \cite{BELL,PUT,HOW} and LP \cite{Dpen} are
classical approaches to MDPs. Ye \cite{Ye,YeSim} gave strongly
polynomial algorithms for a fixed discount factor, including simplex
and policy iteration. Hansen, Miltersen and Zwick \cite{HMZ} extended
the strategy iteration bound to two players with a fixed discount.
Post and Ye \cite{PY} gave a strongly polynomial simplex algorithm
for deterministic MDPs with variable discounts. See also \cite{Uz}.

For a fixed discounted Max MDP, the value vector is the unique minimizer
of the linear program
\[
      \text{minimize }\ \sum_i x_i
      \qquad\text{subject to}\quad x\ge T(x).
\]
There is one inequality for each action. The fixed point is feasible,
and every feasible \(x\) bounds it above by iterating \(x\ge T(x)\).
Theorem~\ref{thm:lp} solves the LP for arbitrary scalar data.

For polynomial data, impose primal feasibility, dual feasibility and
equality of the two objectives, and write the joint system in standard
form \(A(t)y=b(t)\), \(y\ge0\). Let \(m\) be its number of variables
and \(D\) a bound on its coefficient degrees. At each sufficiently
small \(t>0\) there is a basic feasible solution. Along some sequence
tending to zero the same basic columns and independent defining rows
occur. Cramer's rule gives all coordinates over a common denominator,
with numerator and denominator degrees at most \(mD\). The unused
equations hold identically, because they hold on an infinite sequence.
The inequalities hold throughout a sufficiently small positive
interval, because every rational function has a stable sign there.
For related asymptotic LP methods see \cite{Jero,FAA,Vein}.

It suffices, therefore, to seek polynomials \(Y,s\) of degrees at most
\(mD\) satisfying
\[
       A(t)Y(t)=s(t)b(t),\qquad Y(t)\ge0,\qquad s(t)>0
       \quad\text{for all sufficiently small }t>0 .
\]
The equality is a system of linear equations in the coefficients.
For each polynomial in a sign condition, the first nonzero coefficient
in order of increasing degree determines the sign. Starting at degree
zero, impose nonnegativity of the current coefficient in every active
condition. Use Theorem~\ref{thm:lp} to find a feasible point at which
every coefficient that can be positive is positive simultaneously.
The coefficient system is homogeneous, so positivity of a coefficient
can be tested by adding the bound that it be at least one. Average the
witnesses from the successful tests. If a coefficient is forced to zero, prescribe that
equality and advance that condition by one degree. A weak condition
whose coefficients are all zero is satisfied. The strictly positive
denominator condition cannot be exhausted, since each relaxation
contains every true solution. When no coefficient is forced to zero,
the current point satisfies every sign condition. The number of
advances is at most the number of coefficients. Each LP uses the
original coefficient data and the prescribed zeros.

Divide \(Y\) by \(s\) and recover the value coordinates from the
standard-form variables. Choose an action tight identically in \(t\)
at each state. Such an action exists since the action scores are
rational functions and their order stabilizes for small \(t>0\).
This yields a Blackwell policy for~\eqref{eq:gz}, including arbitrary
positive real weights. All degrees are polynomial and the coefficient
arithmetic uses a constant number of layers of sums and products, as
in Section~\ref{sec:arithmetic}.

\subsection{Two players with arbitrary discounts}
\label{sec:longdiscounts}

For two players, the Bellman operator may contract by an arbitrarily
small amount even with a fixed number of states. We will halve the gap
between two policy guarantees by estimating the values on a scale set
by that gap. Only the transition probabilities will enter the slopes
of the fixed-point computation at this scale.

First fix the discounts and rewards. An action \(a\) at state \(i\)
has expected reward \(r_a=\sum_{e\in E_a} p_e r_e\), transition row
\(P_a\), and discounted transition row \(B_a\), where
\[
   P_a x=\sum_{e\in E_a}p_e x_{j(e)},\qquad
   B_a x=\sum_{e\in E_a}p_e\beta_e x_{j(e)} .
\]
Thus the action score in~\eqref{eq:T} is \(r_a+B_a x\). We have
\(0\le B_a\le P_a\) coordinatewise. Write
\(\lambda_a=1-B_a\mathbf1>0\) for the mass lost to discounting.
The row \(P_a-B_a\) has mass \(\lambda_a\).

\paragraph{Policy guarantees.}
Start with any Max policy and any Min policy. Fix Max's policy and
solve the resulting Min MDP to get its vector of guarantees \(L\).
Fix Min's policy and solve the Max MDP to get \(U\). Then
\(L\le v\le U\). If \(W=\|U-L\|_\infty=0\), these policies are
optimal and \(v=L\). For \(W>0\) we construct two new policies whose
width is at most \(W/2\).

\paragraph{The clipped update.}
Given \(y\in\RR^n\) with \(\|v-y\|_\infty\le 2H\), \(H>0\), we seek
a correction \(u\) that halves the error bound. Keep the discounted
score at the center \(y\), propagate the correction using \(P_a\),
and clip each coordinate to \([-H,H]\). Thus we solve
\begin{equation}\label{eq:clip}
 u_i=\left[\mathop{\mathrm{ext}_i}_{a\text{ at }i}
              \Phi_a(u)\right]_{-H}^{H},\qquad
 \Phi_a(u)=r_a+B_a y-y_i+P_a u .
\end{equation}
Here brackets denote projection onto the indicated interval. This
map is nonexpansive, because \(P_a\) is a probability row and clipping
preserves the Lipschitz bound. The discounts enter only the offsets.
After scaling coordinates by \(1/H\), Proposition~\ref{prop:asym}
finds a fixed point in the unit cube. For earlier uses of clipping in
mean-payoff games see \cite{OAIRa,OAIDe}.

Every solution \(u\) of~\eqref{eq:clip} satisfies
\begin{equation}\label{eq:dom}
    \|v-(y+u)\|_\infty\le H .
\end{equation}
To prove the lower bound, translate by the clipping radius and put
\(z=y+u-H\mathbf1\). For every action \(a\) at \(i\),
\begin{equation}\label{eq:compensate}
 r_a+B_a z-z_i
 =\Phi_a(u)-u_i+H\lambda_a-(P_a-B_a)u
 \ \ge\ \Phi_a(u)-u_i ,
\end{equation}
because \(u\le H\mathbf1\) and \(P_a-B_a\) has mass \(\lambda_a\).
If \(u_i>-H\), taking the extremum in~\eqref{eq:compensate} gives
\(z_i\le T_i(z)\). If \(u_i=-H\), the initial error bound gives
\(z_i=y_i-2H\le v_i\). A positive maximum
\(d=\max_i(z_i-v_i)\) would therefore be attained at a state where
\(z_i\le T_i(z)\). But monotonicity and contraction would give
\[
  d\le T_i(z)-T_i(v)
    \le T_i(v+d\mathbf1)-T_i(v)\le\bar\beta d<d .
\]
Thus \(z\le v\). Negate values and exchange the owners to get
\(v\le y+u+H\mathbf1\), proving~\eqref{eq:dom}.

If \(y<v\) componentwise,~\eqref{eq:dom} also forces \(u_i>-H\) at
every state. The maximizing actions for \(\Phi_a(u)\) in~\eqref{eq:clip}
then certify \(z\le T(z)\) in~\eqref{eq:compensate}. They give a Max
policy guaranteeing \(z\ge v-2H\mathbf1\). Figure~\ref{fig:clip}
shows the translation in a two-state game. With discount \(1/2\),
let Max at \(A\) and Min at \(B\) choose between staying and moving
to the other state. Let their action scores be
\[
 T_A(x)=\max\{3/4+x_A/2,\ x_B/2\},\qquad
 T_B(x)=\min\{1+x_B/2,\ -1/4+x_A/2\}.
\]
The fixed point is \(v=(3/2,1/2)\), with Max staying and Min moving.
At the low estimate in the figure, staying at \(A\) gives the larger
clipped-problem action score and thus certifies the lower guarantee
at both states.

\begin{figure}[htbp]
\centering
\begin{tikzpicture}[x=1cm,y=.80cm,font=\small,
  gridline/.style={black!40,thin},
  interval/.style={lagoon,semithick},
  guarantee/.style={ember,line width=1.6pt}]
\newcommand{\erroraxes}[1]{%
  \fill[lagoon!7] (.08,-1) rectangle (3.86,1);
  \foreach \k in {-2,-1,0,1,2}{
    \draw[gridline] (.08,\k)--(3.86,\k);
    \node[left] at (.05,\k){\(\k\)};
  }
  \node[anchor=south,font=\small\bfseries] at (1.95,2.22){#1};
  \node[below] at (1,-2.05){State \(A\)};
  \node[below] at (2.93,-2.05){State \(B\)};
}
\newcommand{\stateerror}[4]{%
  \draw[interval] (#1,{#2-1})--(#1,{#2+1});
  \draw[interval] ({#1-.16},{#2+1})--({#1+.16},{#2+1});
  \draw[#4] ({#1-.16},{#2-1})--({#1+.16},{#2-1});
  \draw[interval,fill=white] (#1,#2) circle[radius=2.6pt];
  \fill[black] (#1,#3) circle[radius=2pt];
}
\begin{scope}
  \erroraxes{Center \(y=(0,2)\)}
  \stateerror{1}{1}{1.5}{interval}
  \stateerror{2.93}{-1}{-1.5}{interval}
\end{scope}
\begin{scope}[xshift=6.3cm]
  \erroraxes{Low center \(y=(0,1/4)\)}
  \stateerror{1}{1}{1.5}{guarantee}
  \stateerror{2.93}{.5}{.25}{guarantee}
\end{scope}
\end{tikzpicture}
\caption{Clipped updates for the two-state game with \(H=1\).
The shaded band is the allowed range of the correction \(u\).
Each vertical interval is \([u_i-1,u_i+1]\), with the correction
at its open circle and the true error \(v_i-y_i\) at its solid dot.
In the right panel \(y<v\) in both states, so the lower endpoints
are certified policy guarantees relative to \(y\). Adding \(y\)
gives the absolute lower bound \(z=(0,-1/4)\).}
\label{fig:clip}
\end{figure}


\paragraph{Halving the width.}
The midpoint \(y=(L+U)/2\) has error at most \(W/2\). Apply the
clipped update three times with radii \(W/4,W/8,W/16\), replacing
\(y\) by \(y+u\) each time. By~\eqref{eq:dom} the resulting
\(y_*\) has error at most \(W/16\). Set \(H=W/8\). The lower estimate
\(y_*-H\mathbf1\) is strictly below \(v\) with error at most \(2H\),
so one more clipped solve gives a Max policy guaranteeing at least
\(v-(W/4)\mathbf1\). Likewise \(y_*+H\mathbf1\) gives a Min policy
guaranteeing at most \(v+(W/4)\mathbf1\). The two new MDP guarantees
thus have width at most \(W/2\).

\paragraph{The possible scales of a policy error.}
Halving alone gives an approximation. To get exact termination in
polynomially many iterations, we show that each positive width lies
within a factor \(2^{p(N)}\) of one of polynomially many positive
numbers, for a fixed polynomial \(p\). Here \(N\) is the length of the
discrete input, including the probability bits. The scales may depend
on the real input. The algorithm does not compute them.

Each of \(L,U\) is the payoff vector of a pair of policies, since a
best response in an MDP is positional and works from every state.
Moreover,
\[
   \max\{\|v-L\|_\infty,\|U-v\|_\infty\}\le W
      \le 2\max\{\|v-L\|_\infty,\|U-v\|_\infty\}.
\]
It therefore suffices to control the error of each policy pair. Write
\(\Delta\) for the vector of its payoffs minus \(v\) and suppose that
\(G=\|\Delta\|_\infty>0\). The \emph{action gaps}
\(\delta_a=r_a+B_a v-v_i\), defined for every action \(a\) at \(i\),
are common to all pairs. For the action used by the pair at \(i\),
\[
     \Delta_i=\delta_a+B_a\Delta .
\]
Thus \(g=\max_{a\ {\rm used}}|\delta_a|\le 2G\), and \(g>0\) by
contraction. When \(G\) is much larger than \(g\), the transition
chain must retain a large error across a whole closed set of states.

Indeed, change the signs of \(\Delta,\delta\) if necessary so that
some \(\Delta_i=G\). Write \(D_i=G-\Delta_i\). For the used action
\(a\) at \(i\),
\[
   \lambda_aG+B_aD=D_i+\delta_a\le D_i+g,\qquad
   P_aD\le 2(D_i+g).
\]
For the second inequality use \(D\le2G\mathbf1\) and
\((P_a-B_a)\mathbf1=\lambda_a\). Every positive transition has
probability at least \(2^{-\mathrm{poly}(N)}\). Along a transition
\(i\to j\), the last display bounds \(D_j\) by
\(2^{\mathrm{poly}(N)}(D_i+g)\). Each state reachable from a maximum
can be reached by a path of length at most \(n-1\). Starting with
\(D_i=0\) at the maximum, we obtain \(D_j\le Kg\) throughout its
reachable set for some uniform \(K=2^{\mathrm{poly}(N)}\).

If \(G\le2Kg\), the single gap of magnitude \(g\) supplies the scale.
Otherwise \(G/2\le\Delta_j\le G\) on the reachable set. This set is
closed under the undiscounted transitions of the pair. Choose a
stationary distribution supported there, assigning mass \(d_a\) to the
used action in each state. Such a distribution can be taken at a vertex
of the stationarity equations with nonnegative coordinates summing to
one. These equations use only transition probabilities, so every
positive \(d_a\) has polynomial bit length. Stationarity and the
recursion for \(\Delta\) give
\[
    \sum_a d_a\delta_a=\sum_a d_a(P_a-B_a)\Delta .
\]
The rows \(P_a-B_a\) are nonnegative, stay in the same reachable set,
and have masses \(\lambda_a\). Hence
\[
   (G/2)\sum_a d_a\lambda_a
   \le \sum_a d_a\delta_a\le G\sum_a d_a\lambda_a .
\]
The weighted discount loss \(\sum_a d_a\lambda_a\) is within a factor
\(2^{\mathrm{poly}(N)}\) of the largest \(\lambda_a\) in the support.
The numerator in the middle is a nonzero short rational combination
of the action gaps.

It remains to bound the number of scales for such numerators. Fix the
polynomial bit bound for their coefficient vectors. Take a maximal
collection of nonzero magnitudes of combinations of the gaps that
differ pairwise by factors greater than \(2^{p(N)}\), where \(p\) is
chosen below. By maximality each nonzero magnitude is within that
factor of one in the collection. The selected coefficient vectors are
linearly independent. Indeed, any dependence contains a dependence on
at most \(m+1\) vectors, where \(m\) is the number of actions. Clearing
denominators and applying determinant bounds gives integer dependence
coefficients of \(\ell_1\)-norm less than \(2^{p(N)}\) for a
sufficiently large fixed polynomial \(p\). Applying that dependence to
the combination values gives a contradiction. The term of largest
magnitude among the values with nonzero coefficient exceeds the sum of
the magnitudes of all smaller terms even after weighting by their
coefficients.

There are at most \(m\) numerator scales. For \(G>2Kg\), divide one
of them by one of the \(m\) discount losses \(\lambda_a\). For
\(G\le2Kg\), use one of the \(m\) individual gap magnitudes. After
increasing \(p\) to absorb the comparison factors, these at most
\(m^2+m\) scales work for every pair, hence for every positive \(W\).
A sequence of successive halvings has at most \(2p(N)+1\) widths
in the interval from \(2^{-p(N)}\) to \(2^{p(N)}\) times any one scale.
The algorithm therefore reaches \(W=0\) after polynomially many
iterations.

\paragraph{Arithmetic and parameter dependence.}
Only the two discrete policies pass between iterations. Within each
iteration a constant number of MDP evaluations and clipped solves
suffice. The step count is polynomial in \(N\), and the bit bounds of
Section~\ref{sec:arithmetic} apply to each iteration.

Finally let the rewards and discounts be polynomials in \(t\), valid
for all sufficiently small positive \(t\). Their degrees count in the
input size, and the probabilities are fixed binary rationals. The MDP
algorithm of Section~\ref{sec:mdp} returns best responses and
guarantees as rational functions. Compare two such functions using
their signs at \(0^+\), so the policy selections hold throughout a
sufficiently small positive interval.

In a clipped solve the height \(H\) and offsets
\(c_a=r_a+B_a y-y_i\) are now rational functions. Choose a common
denominator \(d\), positive near \(0^+\), so that \(dH\) and every
\(dc_a\) are polynomials. Let \(K>0\) bound \(|dH|\) on \(0<t<1\),
for example one plus the sum of the absolute values of its
coefficients. In the scaled coordinates \(\xi=du/K\), Equation
\eqref{eq:clip} has offsets \(dc_a/K\), rows \(P_a\), and clipping
endpoints \(\pm dH/K\). These endpoints lie in \([-1,1]\) near
\(0^+\). The resulting nonexpansive map of the unit cube fits
Proposition~\ref{prop:asym}, which yields a rational \(\xi\) and the
required extremizing actions. Recover \(u=K\xi/d\).

Degrees remain polynomial because each iteration has a constant number
of layers of calls, and passes only policies to the next iteration.
Products and sums of polynomially many polynomials of these degrees
can be computed by evaluation and interpolation using polynomial-size
formulas in their coefficients, as required in
Section~\ref{sec:arithmetic}. Every finite sequence of symbolic
iterations agrees with a run at a common sufficiently small \(t>0\).
The pointwise iteration bound is uniform in \(t\), so the symbolic
algorithm has the same bound. Applied to~\eqref{eq:gz}, it supplies
Blackwell optimal policies even for arbitrary positive real weights.
The mean-payoff limit of Section~\ref{sec:shortslopes} completes the
proof of Theorem~\ref{thm:games}.
